\chapter{Impl\'ementation ESP32}

Ce chapitre d\'eveloppe en d\'etail l'impl\'ementation du firmware de l'ESP32 pour le n\oe{}ud Andon, en pr\'esentant le code source, les biblioth\`eques utilis\'ees et les m\'ecanismes de communication.

\section{Environnement de d\'eveloppement}

\subsection{Arduino IDE}

Le firmware ESP32 est d\'evelopp\'e avec Arduino IDE, qui offre :

\begin{itemize}
    \item Un environnement de d\'eveloppement intuitif.
    \item La compilation et le t\'el\'echargement direct vers l'ESP32.
    \item Un large \'ecosyst\`eme de biblioth\`eques.
    \item Le d\'ebugage s\'erial pour le suivi d'ex\'e-cution.
\end{itemize}

\subsection{Biblioth\`eques utilis\'ees}

\begin{table}[htbp]
\centering
\renewcommand{\arraystretch}{1.3}
\begin{tabular}{l l}
\toprule
\textbf{Biblioth\`eque} & \textbf{Fonction} \\
\midrule
WiFi.h & Connexion Wi-Fi \\
PubSubClient & Client MQTT \\
ArduinoJson & G\'en\'eration JSON \\
WiFiClientSecure & Connexion TLS \\
\bottomrule
\end{tabular}
\caption{Biblioth\`eques Arduino utilis\'ees}
\label{tab:bibliotheques_arduino}
\end{table}

\section{Code source du firmware}

\subsection{Configuration initiale}

\begin{lstlisting}[language=C++,
caption={Configuration et constantes},
label={lst:config_esp32}]
#include <WiFi.h>
#include <PubSubClient.h>
#include <ArduinoJson.h>

// Configuration WiFi
const char* SSID = "SEWS_IOT";
const char* PASSWORD = "****";

// Configuration MQTT
const char* MQTT_SERVER = "192.168.1.100";
const int MQTT_PORT = 1883;
const char* MQTT_USER = "esp32_ligne1";
const char* MQTT_PASS = "****";

// Configuration ligne
const int LIGNE_ID = 1;
const char* TOPIC_ALERT = 
  "sews/andon/1/alert";
const char* TOPIC_LED = 
  "sews/command/1/led";

// Broches GPIO
const int BOUTON_PIN = 4;
const int LED_VERTE = 16;
const int LED_ROUGE = 17;

WiFiClient wifiClient;
PubSubClient mqtt(wifiClient);

unsigned long dernierAppui = 0;
const unsigned long DEBOUNCE = 200;
\end{lstlisting}

\subsection{Connexion WiFi et MQTT}

\begin{lstlisting}[language=C++,
caption={Connexion WiFi et MQTT},
label={lst:connect_esp32}]
void connectWiFi() {
  Serial.print("Connexion WiFi...");
  WiFi.begin(SSID, PASSWORD);
  
  while (WiFi.status() != WL_CONNECTED) {
    delay(500);
    Serial.print(".");
  }
  
  Serial.println(" OK");
  Serial.println(WiFi.localIP());
}

void connectMQTT() {
  mqtt.setServer(MQTT_SERVER, MQTT_PORT);
  mqtt.setCallback(callbackMQTT);
  
  while (!mqtt.connected()) {
    Serial.print("Connexion MQTT...");
    
    if (mqtt.connect(
        "esp32_ligne1", MQTT_USER, MQTT_PASS
    )) {
      Serial.println(" OK");
      mqtt.subscribe(TOPIC_LED);
    } else {
      Serial.println(" echec");
      delay(5000);
    }
  }
}
\end{lstlisting}

\subsection{Boucle principale}

\begin{lstlisting}[language=C++,
caption={Boucle principale (loop)},
label={lst:loop_esp32}]
void setup() {
  Serial.begin(115200);
  
  pinMode(BOUTON_PIN, INPUT);
  pinMode(LED_VERTE, OUTPUT);
  pinMode(LED_ROUGE, OUTPUT);
  
  digitalWrite(LED_VERTE, HIGH);
  digitalWrite(LED_ROUGE, LOW);
  
  connectWiFi();
  connectMQTT();
}

void loop() {
  if (!mqtt.connected()) {
    connectMQTT();
  }
  mqtt.loop();
  
  // Detection appui bouton
  if (digitalRead(BOUTON_PIN) == HIGH) {
    unsigned long now = millis();
    if (now - dernierAppui > DEBOUNCE) {
      dernierAppui = now;
      envoyerAlerte();
      
      // Signalisation locale
      digitalWrite(LED_VERTE, LOW);
      digitalWrite(LED_ROUGE, HIGH);
    }
  }
}
\end{lstlisting}

\subsection{Envoi d'alerte Andon}

\begin{lstlisting}[language=C++,
caption={Fonction d'envoi d'alerte},
label={lst:envoyer_alerte_esp32}]
void envoyerAlerte() {
  // Construction du message JSON
  StaticJsonDocument<256> doc;
  doc["ligneId"] = LIGNE_ID;
  doc["timestamp"] = millis();
  doc["type"] = "panne";
  doc["urgence"] = "haute";
  doc["ip"] = WiFi.localIP().toString();
  
  char buffer[256];
  serializeJson(doc, buffer);
  
  // Publication MQTT
  bool success = mqtt.publish(
    TOPIC_ALERT, buffer, true
  );
  
  if (success) {
    Serial.println("Alerte envoyee");
  } else {
    Serial.println("Echec envoi");
  }
}
\end{lstlisting}

\subsection{R\'e-ception de commandes}

\begin{lstlisting}[language=C++,
caption={Callback MQTT pour les commandes LED},
label={lst:callback_mqtt}]
void callbackMQTT(
  char* topic, byte* payload, unsigned int len
) {
  // Parser le message JSON
  StaticJsonDocument<128> doc;
  deserializeJson(doc, payload, len);
  
  bool rouge = doc["rouge"];
  bool verte = doc["verte"];
  
  digitalWrite(LED_ROUGE, rouge ? HIGH : LOW);
  digitalWrite(LED_VERTE, verte ? HIGH : LOW);
  
  Serial.println("LEDs mises a jour");
}
\end{lstlisting}

\section{Gestion de la fiabilit\'e}

\subsection{Reconnexion automatique}

Le firmware int\`egre un m\'ecanisme de reconnexion automatique :

\begin{lstlisting}[language=C++,
caption={Reconnexion MQTT},label={lst:reconnect}]
void reconnectMQTT() {
  while (!mqtt.connected()) {
    Serial.print("Reconnexion MQTT...");
    
    if (mqtt.connect(
        "esp32_ligne1", MQTT_USER, MQTT_PASS
    )) {
      Serial.println("OK");
      mqtt.subscribe(TOPIC_LED);
      mqtt.publish(
        "sews/capteurs/1/status",
        "{\"status\":\"online\"}", true
      );
    } else {
      Serial.print("Echec, rc=");
      Serial.print(mqtt.state());
      Serial.println(" Retry 5s");
      delay(5000);
    }
  }
}
\end{lstlisting}

\subsection{Surveillance de la connexion WiFi}

\begin{lstlisting}[language=C++,
caption={Surveillance WiFi},label={lst:watchdog_wifi}]
void verifierWiFi() {
  if (WiFi.status() != WL_CONNECTED) {
    Serial.println("WiFi perdu...");
    WiFi.reconnect();
    
    int tentatives = 0;
    while (WiFi.status() != WL_CONNECTED 
           && tentatives < 20) {
      delay(500);
      tentatives++;
    }
    
    if (WiFi.status() == WL_CONNECTED) {
      Serial.println("WiFi restaure");
    }
  }
}
\end{lstlisting}

\section{Tests du firmware}

\subsection{Tests unitaires}

Le firmware a \'et\'e test\'e sur les aspects suivants :

\begin{itemize}
    \item D\'etection d'appui sur le bouton avec d\'ebounce.
    \item Allumage/extinction des LED.
    \item Connexion WiFi et MQTT.
    \item Publication et r\'e-ception de messages.
    \item Reconnexion apr\`es d\'econnexion.
\end{itemize}

\subsection{Tests en conditions r\'eelles}

Le firmware a \'et\'e test\'e sur l'atelier Test \'Electrique pendant 2 semaines :

\begin{itemize}
    \item Stabilit\'e de la connexion Wi-Fi ($>$ 99,5\% de disponibilit\'e).
    \item Temps de r\'eponse moyen de l'alerte : 180 ms.
    \item Consommation : 80 mA en fonctionnement normal.
    \item R\'e-sistance aux perturbations \'electromagn\'etiques.
\end{itemize}

\section{Synth\`ese du chapitre}

Le firmware ESP32 du n\oe{}ud Andon est impl\'ement\'e en Arduino/C++ avec les biblioth\`eques WiFi, PubSubClient et ArduinoJson. Il assure la d\'etection des appuis, l'envoi d'alertes MQTT, la r\'e-ception de commandes LED et la reconnexion automatique. Les tests ont d\'emontr\'e la fiabilit\'e du syst\`eme en conditions industrielles r\'eelles.
